Op de oorspronkelijke beelden herstelt SVM 157 fouten en verliest het 293 goede voorspellingen ten opzichte van niets aanpassen; netto -136. De foldverschillen zijn -51, -53, -28, -3, -1. Tegenover spiegeltraining is het netto verschil -110. Normalisatie verbetert de oorspronkelijke classificatie dus niet.

Het aparte mengsel heeft 60.225 geldige detecties. Daar verandert SVM-macro-F1 van 0,62650 naar 0,96047; spiegeltraining geeft 0,97421. Dit is een andere toetsverdeling, niet de oorspronkelijke baseline. Het uitsluiten van detectiemislukkingen verandert hier de precisie niet: die pogingen leveren geen voorspelde letter en dus geen false positive. F1 verandert wel doordat hun false negatives wegvallen. Alle-pogingenscores blijven apart in de resultaatbestanden staan.
